% SGA 6: índice terminológico español alineado con la fuente.
% Alcance: líneas 15923--16062 inclusive de la fuente francesa.
% Fuente francesa: sga6_fr_r1/src/fr.tex.
% SHA-256 de la fuente francesa completa:
% 0D2158B9FE3B466309A92165727A0B191422E05FBAA0DACD5806E8322C54BA1A.

\section*{Índice terminológico}
\begingroup\small
\begin{longtable}{@{}>{\raggedright\arraybackslash}p{0.64\linewidth}>{\raggedright\arraybackslash}p{0.14\linewidth}>{\raggedright\arraybackslash}p{0.12\linewidth}@{}}
\toprule
 & Exp. & §\\
\midrule
\endhead
aditiva (aplicación) & I & 6.1\\
aditiva (aplicación) & IV & 1.5\\
\(K-\lambda\)-álgebra aumentada & V & 3.9\\
algebraicamente equivalente a cero (en \(K^\bullet\)) & X App. & 7.11\\
algebraicamente equivalente a cero (haz invertible) & XIII & 4.4\\
de amplitud \(\mathcal C_0\)-perfecta contenida en \([a,b]\) (complejo) & I & 4.7\\
de amplitud \(\mathcal C_0\)-perfecta finita (complejo) & I & 4.7\\
de amplitud perfecta contenida en \([a,b]\) (complejo) & I & 4.8\\
de amplitud perfecta finita (complejo) & I & 4.8\\
de amplitud plana contenida en \([a,b]\) (complejo) & I & 5.2\\
de amplitud plana finita (complejo) & I & 5.2\\
de amplitud perfecta relativa a \(f\) (o a \(Y\)) contenida en \([a,b]\) (complejo) & III & 4.1\\
de amplitud perfecta contenida en \([-n,0]\) (morfismo) & III & 4.1\\
de amplitud perfecta finita relativa a \(f\) (complejo) & III & 4.1\\
de amplitud plana relativa a \(f\) (o a \(Y\)) contenida en \([a,b]\) (complejo) & III & 3.1\\
de amplitud plana contenida en \([-n,0]\) (morfismo) & III & 3.1\\
de amplitud plana finita relativa a \(f\) (complejo) & III & 3.1\\
amplitud proyectiva & I & 4.12.1\\
\(\lambda\)-anillo & V & 2.4\\
\(\lambda\)-anillo engendrado por generadores y relaciones & V & 4.6\\
\(\lambda\)-anillo libre engendrado por una familia de generadores & V & 4.4\\
anillo de Chern & V & 6.2\\
aplicación polinómica & X & 2.1.1\\
binomial (anillo) & V & 2.7\\
carácter de Chern & V & 6.4\\
carácter de Todd & V & 1.16\\
\(S\)-categoría abeliana & I & 1.1.1\\
\(S\)-categoría abeliana plana & I & 1.1.1\\
\(S\)-categoría aditiva & I & 1.1.1\\
\(S\)-categoría triangulada & I & 1.1.1\\
clase de Chern & V & 6.7\\
clase de Chern completada & V & 6.7\\
clase de Chern total & V & 6.7\\
codimensión (de una inmersión regular) & VII & 1.6\\
coherente (complejo) & I & 3.1\\
coherador & II & 3.2\\
complejo cotangente relativo & VIII & 2.3\\
complejo elíptico & II App. II & 1.2.5\\
producto cup & I & 8.2\\
\(k\)-ciclo positivo especial & XIII & 6.9\\
grado (de una aplicación polinómica) & X & 2.1.1\\
grado (de un haz invertible) & XIII & 2.2\\
grado (de un morfismo de esquemas) & X & 4.4\\
denso (subconjunto de una categoría abeliana o triangulada) & IV & 1.7\\
dimensión perfecta & I & 4.14.1\\
dimensión relativa virtual & VIII & 1.9\\
divisorial & II & 2.2.5\\
dual & I & 7.3\\
dualidad perfecta & I & 7.8\\
equivalencia numérica (en \(\operatorname{Gr}^{\bullet}(X)\) y \(\operatorname{Gr}_{\cdot}(X)\)) & X & 4.5\\
\(\tau\)-equivalencia & X & 2.4.1\\
\(\tau\)-equivalencia & XIII & 4.4\\
espacio pseudoanalítico complejo & II & 2.4.7\\
espacio pseudoanalítico complejo de Stein & II & 2.4.7\\
exacto (subconjunto de una categoría abeliana o triangulada) & IV & 1.7\\
haz \((b)\) & XIII & 1.5\\
haz conormal & VII & 1.6\\
haz cotangente relativo & II App. II & 1.2.1\\
haz tangente relativo & II App. II & 1.2.1\\
familia \((b)\) & XIII & 1.12\\
familia amplia & II & 2.2.4\\
familia de clases de haces coherentes & XIII & 1.12\\
familia acotada de clases de haces coherentes & XIII & 1.12\\
familia de clases de haces coherentes acotada por un haz coherente & XIII & 1.12\\
fibrado cotangente relativo & II App. II & 1.2.1\\
fibrado tangente relativo & II App. II & 1.2.1\\
fielmente plano (morfismo de topoi anillados) & I & 2.16.2\\
filtración de una \(\lambda\)-álgebra aumentada & V & 3.10\\
\(S\)-funtor aditivo & I & 1.1.3\\
\(S\)-funtor exacto & I & 1.1.3\\
fórmula de proyección & IV & 2.11\\
grupo de Grothendieck & IV & 2.2\\
\(\lambda\)-homomorfismo & V & 2.1\\
homomorfismo de especialización & X App. & 7.3\\
homomorfismo de especialización & X App. & 7.16\\
\(\lambda\)-ideal & V & 4.5\\
índice analítico & II App. II & 3.2.2\\
de intersección completa (morfismo) & VIII & 1.1\\
\(n\)-isomorfismo & I & 1.4.3\\
libre homogéneo (módulo graduado) & VI & 1.2.2\\
localmente estable por núcleo de epimorfismo & I & 1.2\\
localmente cuasilevantable & I & 1.2\\
localmente cuasiestable por núcleo de epimorfismo & I & 1.2\\
localmente levantable & I & 1.2\\
número de Chern & X & 4.2\\
número de intersección & X & 4.1\\
número de intersección (para una familia de haces invertibles) & X & 4.4\\
numéricamente equivalente a cero (haz invertible) & XIII & 4.4\\
operador diferencial universal de orden \(n\) & II App. II & 1.2.1\\
operaciones de Adams & V & 7.1\\
\(\mathcal C_0\)-perfecto (complejo) & I & 4.7\\
perfecto (complejo) & I & 4.8\\
perfecto (morfismo) & III & 4.1\\
perfecto relativo a \(f\) (complejo) & III & 4.1\\
precoherente & I & 3.1\\
de \(n\)-\(\mathcal C_0\)-presentación finita & I & 2.8\\
de presentación finita & I & 2.8\\
de \(\infty\)-\(\mathcal C_0\)-presentación finita & I & 2.8\\
\(n\)-\(\mathcal C_0\)-pseudocoherente (complejo) & I & 2.2\\
\(\mathcal C_0\)-pseudocoherente (complejo) & I & 2.2\\
\(n\)-pseudocoherente (complejo) & I & 2.15\\
\(n\)-pseudocoherente (morfismo) & III & 1.2\\
pseudocoherente (complejo) & I & 2.15\\
pseudocoherente (morfismo) & III & 1.2\\
\(n\)-pseudocoherente relativo a \(f\) (complejo) & III & 1.2\\
pseudocoherente relativo a \(f\) (complejo) & III & 1.2\\
polinomio \((b)\) & XIII & 1.9\\
polinomio de Snapper & X & 2.3.1\\
polinomio universal & V & 1.15\\
pre-\(\lambda\)-anillo & V & 2.1\\
\(n\)-cuasiisomorfismo & I & 1.4.3\\
rango (de un complejo perfecto) & I & 6.11\\
\(m\)-regular (haz) & XIII & 1.1\\
regular (homomorfismo) & VII & 1.1\\
regular (ideal) & VII & 1.4\\
regular (inmersión) & VII & 1.4\\
sub-\(S\)-categoría triangulada & I & 1.1.4\\
especial (\(\lambda\)-anillo) & V & 2.4\\
especial (morfismo) & III & 1.1.4\\
de Stein (topos localmente anillado) & II & 2.4.1\\
estrictamente \(\mathcal C_0\)-perfecto & I & 2.1\\
estrictamente \(\mathcal C_0\)-\(n\)-pseudocoherente & I & 2.1\\
estrictamente \(\mathcal C_0\)-pseudocoherente & I & 2.1\\
sistema mínimo de generadores & VII & 1.4.1\\
de dimensión Tor finita & I & 5.2\\
independencia Tor & III & 1.5\\
traza & I & 8.1\\
de \(\mathcal C_0\)-tipo finito & I & 1.2\\
variedad mixta de clase \(C^r\) & II App. II & 1.1\\
\bottomrule
\end{longtable}
\endgroup
\newpage
